Torricelli Asymptotic Cusp Encryption (TACE) v3:
A Non-Archimedean Path-Commitment Primitive
and Its Algebraic Funnel Analysis
Yusuf Arman Acin
September 17, 2026
Abstract
This paper introduces Torricelli Asymptotic Cusp Encryption (TACE) v3, a novel crypto-
graphic path-commitment primitive built upon non-Archimedean p-adic ultrametric trees.
Motivated geometrically by the asymptotic properties of Gabriel’s Horn—where an infi-
nite surface area encapsulates a finite volume—TACE maps hierarchical combinatorial path
spaces into cryptographic commitments. Through empirical cryptanalysis and backward-
branching simulation, we demonstrate that the core modular recurrence functions as an al-
gebraic “funnel” with tightly bounded predecessor trajectories rather than an unconstrained
chaotic expansion. To secure this structure against multi-path collision vectors, TACE v3
integrates a domain-separated cryptographic hash chain that cryptographically binds the
entire state-digit traversal. Finally, we establish a formal quantum security analysis demon-
strating that the absence of a global abelian hidden-subgroup structure prevents the direct
application of Shor’s Quantum Fourier Transform (QFT) period-finding decimation.
1 Introduction and Geometric Motivation
Traditional public-key cryptosystems rely heavily on algebraic problems defined over finite fields
or cyclic groups, rendering them fundamentally vulnerable to quantum algorithms that exploit
cyclic group periodicities via the Quantum Fourier Transform (QFT).
TACE explores an alternative topological foundation by abandoning flat Euclidean grids
in favor of non-Archimedean ultrametric spaces. Geometrically, this is inspired by Torricelli’s
Paradox (Gabriel’s Horn), where an infinite surface area is associated with a finite volume.
In TACE, this intuition is translated into discrete digital architecture: an exponentially large
combinatorial path space is compressed into a compact cryptographic commitment anchor.
Rather than presenting an asymmetric encryption scheme with an algebraic trapdoor, TACE
v3 is formally defined as a **Non-Archimedean Path-Commitment Primitive** (or path-hiding
commitment function), providing rigorous guarantees for path integrity and non-malleability.
2 Mathematical Framework: Ultrametric p-Adic Trees and Com-
mitment
TACE operates over the ring of p-adic integers Zp, structured through hierarchical tree traver-
sals.
2.1 The Private Path Space
The private key is a finite sequence of digits
k = (a0, a1, . . . , an−1),
1
where each ai ∈ {0, 1, . . . , p − 1}, representing a unique path down an N -level hierarchical tree
of total theoretical size |P| = pN .
2.2 The Cusp Transition Function
Successive nodes evolve via a nonlinear polynomial lifting inspired by Hensel’s Lemma:
xi+1 = (xi · p + ai)2 + i (mod pi+1), x0 = 0.
This recurrence governs the forward evolution of the state through expanding prime-power
moduli while preserving the non-Archimedean hierarchical structure.
2.3 The Path-Commitment Chain (Public Anchor)
Empirical inversion testing reveals that the raw recurrence is lossy and many-to-one, forming
an algebraic funnel rather than a strict permutation. To prevent multi-key collision forgery,
TACE v3 binds the entire trajectory using a domain-separated cryptographic hash chain:
c0 = H(“TACE-INIT” ∥ x0)
ci+1 = H(“TACE-STEP” ∥ i ∥ ci ∥ xi ∥ ai), for i = 0, . . . , N − 1.
The definitive Public Anchor is the final accumulator state:
Ck = cN
Even if two distinct private key sequences share a terminal residue due to modular folding,
their intermediate states or chosen digits diverge, forcing the hash chain commitment Ck to be
entirely unique.
3 Empirical Cryptanalysis: The Algebraic Funnel Hypothesis
To evaluate the internal behavior of the recurrence, automated backward-branching and prede-
cessor distribution experiments were conducted across various parameter sets (e.g., p = 5, N =
8).
3.1 Backward Branching Trajectory
When tracing backward from a target terminal state xN down to Level 0, the active predecessor
states do not explode exponentially, nor do they immediately collapse to zero. Instead, the
quadratic congruence constraints act as a narrow algebraic sieve:
• Across tested state spaces, the backward search fringe remains tightly bounded (e.g.,
peaking at fewer than 30 active states out of hundreds of thousands of possibilities).
• This empirical finding proves that the raw recurrence alone is insufficient as a standalone
cryptographic key without a commitment wrapper, as narrow backward paths could the-
oretically be traversed if unconstrained.
3.2 Vindication of the Hash Chain Architecture
The discovery of this algebraic funnel confirms that TACE’s security must rely on **path-
binding**. While the modular recurrence provides the ultrametric hierarchical structure, the
hash chain Ck = cN ensures that an adversary cannot forge a valid interaction using a colliding
alternative path.
2
4 Quantum Security Analysis
Theorem 1 (Absence of Abelian Hidden-Subgroup Structure in TACE). Let TN be a truncated
p-adic hierarchical tree of depth N over Zp governed by the nonlinear transition function fi(x) =
(x· p+ ai)2 +i (mod pi+1). The path space P = {0, . . . , p−1}N mapped via recursive polynomial
evaluation admits no global compact abelian group structure with a translation-invariant cyclic
generator. Consequently, the standard abelian hidden-subgroup formulation of Shor’s algorithm
via the Quantum Fourier Transform (QFT) does not directly apply to TACE’s path-recovery
problem.
Proof. 1. Shor’s algorithm and its period-finding subroutine rely fundamentally on evaluating a
periodic function f (x) = f (x + r) over a finite abelian group G = Zm. 2. The TACE recurrence
introduces level-dependent shifts +i and quadratic cross-terms within expanding prime power
moduli pi+1, rendering the mapping non-homomorphic with respect to standard group addition.
3. The resulting ultrametric topology (Zp, dp) satisfies the strong triangle inequality and is
totally disconnected, lacking a global translation-invariant offset r. 4. Without a translation-
invariant cyclic subgroup structure, the quantum characters χk(x) = e2πikx/M fail to generate
the orthogonal phase relationships required for QFT period decimation. ■
5 Conclusion and Future Work
Torricelli Asymptotic Cusp Encryption (TACE) v3 establishes a rigorous framework for non-
Archimedean path-commitment primitives. By pairing ultrametric hierarchical tree structures
with a cryptographic hash-chain accumulator, TACE successfully addresses modular folding
collisions while resisting standard QFT period-finding attacks.
Future work will focus on formalizing the Ultrametric Path Inversion (UPI) hardness as-
sumption, executing systematic predecessor distribution scaling across larger primes, and mod-
eling query complexities against structured quantum search baselines.
Data and Code Availability Statement
The simulation harnesses, predecessor distribution analyzers, and path-commitment verification
scripts supporting this research are available upon request for academic peer-review
